\chapter{共形对称与最高权表示}
\label{chap:ghs-conformal-representation-ope}

Poincar\'e 对称把单粒子态按质量与自旋分类，但它不把不同长度尺度联系起来。
若一个理论的局域结构在尺度变换下保持不变，平移后的算符与它的各阶导数便受到额外约束；
这时质量壳不再是最自然的表示标签，标度维数和旋转表示取而代之。
前面用最高权与升降算符构造的思想仍然有效，不过现在平移生成全部后裔，特殊共形生成元反向降低后裔级数。

对称无迹张量的去迹与旋转分解已经在\secref{sec:ghs-harmonics-stf}中建立。
此处让这些旋转表示带上标度维数，研究正定内积会允许哪些维数，以及哪些导数方向必须在表示中模掉。
这些是共形代数的结果；具体自由场的双线性构造、径向反射正性和相关函数展开可在《自由电子量子光学讲义》的径向量子化章中继续使用。

\section{共形变换与生成元}
\label{sec:ghs-conformal-algebra-unitarity}

\begin{definition}[欧氏共形变换]
以固定参考长度将坐标无量纲化。在 $d>2$ 维欧氏空间，一个局部可逆坐标变换 $x\mapsto x'$ 称为共形变换，若
\[
 \delta_{\alpha\beta}\frac{\partial x'^\alpha}{\partial x^\mu}
 \frac{\partial x'^\beta}{\partial x^\nu}=\Omega(x)^2\delta_{\mu\nu},\qquad\Omega(x)>0.
\]
它保持夹角，而允许长度随位置缩放。对 $x'^\mu=x^\mu+\epsilon v^\mu(x)$ 展开一阶，得到共形 Killing 方程
\begin{equation}\label{eq:ghs-conformal-killing}
 \partial_\mu v_\nu+\partial_\nu v_\mu
 =\frac{2}{d}(\partial\cdot v)\delta_{\mu\nu}.
\end{equation}
\end{definition}

\begin{example}[特殊共形变换的局部生成]
把反演、平移与反演复合，求特殊共形变换，展开其无穷小形式并验证共形 Killing 方程。

\textbf{解。} 反演 $x\mapsto x/x^2$，再作平移 $y\mapsto y-b$，最后再次反演，给出
\[
 x'^\mu=\frac{x^\mu-b^\mu x^2}{1-2b\cdot x+b^2x^2}.
\]
这个表达式只在分母不为零的坐标区域使用。
展开至 $b$ 的一阶得到 $v^\mu=2(b\cdot x)x^\mu-b^\mu x^2$，于是
\[
 \partial_\mu v_\nu+\partial_\nu v_\mu=4(b\cdot x)\delta_{\mu\nu},
 \qquad\partial\cdot v=2d(b\cdot x).
\]
两式满足\eqnrefs{eq:ghs-conformal-killing}。平移、旋转、尺度与这一变换共有
$d+d(d-1)/2+1+d=(d+1)(d+2)/2$ 个生成元；其欧氏代数是 $\mathfrak{so}(d+1,1)$。
前章的时向正号 Minkowski 空间仍使用 $\eta$；此处使用欧氏 $\delta$，两种实形式不直接混用。
\end{example}

\begin{definition}[共形生成元对标量的作用]
在 $d>2$ 维欧氏空间，局域共形代数由平移 $P_\mu$、旋转 $M_{\mu\nu}=-M_{\nu\mu}$、尺度变换 $D$ 和特殊共形变换 $K_\mu$ 构成。这里取不含额外虚数单位的径向约定，生成元作用于局域标量场的交换子为
\begin{align}\label{eq:ghs-conformal-generators-scalar}
 [P_\mu,\mathcal O(x)]&=\partial_\mu\mathcal O(x),\\
 [D,\mathcal O(x)]&=(x\cdot\partial+\Delta)\mathcal O(x),\notag\\
 [M_{\mu\nu},\mathcal O(x)]&=(x_\nu\partial_\mu-x_\mu\partial_\nu)\mathcal O(x),\notag\\
 [K_\mu,\mathcal O(x)]&=\bigl[2x_\mu(x\cdot\partial+\Delta)-x^2\partial_\mu\bigr]\mathcal O(x).\notag
\end{align}
\end{definition}

旋转张量的后两式还须附加自旋表示矩阵；在原点，旋转不再作用于坐标，却仍会旋转自由指标。由这些微分算符直接计算得到
\begin{align}\label{eq:ghs-conformal-commutators}
 [D,P_\mu]&=P_\mu,&[D,K_\mu]&=-K_\mu,\\
 [K_\mu,P_\nu]&=2\delta_{\mu\nu}D-2M_{\mu\nu},&
 [M_{\mu\nu},P_\rho]&=\delta_{\nu\rho}P_\mu-\delta_{\mu\rho}P_\nu.\notag
\end{align}
另有通常的 $\mathfrak{so}(d)$ 旋转交换子，以及 $[P_\mu,P_\nu]=[K_\mu,K_\nu]=0$。与前文带 $\ii$ 的厄密角动量生成元相比，这里的 $M_{\mu\nu}$ 是实反对称矩阵约定；两个约定之间只差生成元的虚数因子，不能把它们的交换子直接混写。

在固定坐标 $x$ 上，交换子作用于场算符。计算复合交换子时，显式的坐标系数不是场算符，
所以不能把场上的伴随作用误认为只对坐标函数作用的向量场括号。例如令 $L_D=x\cdot\partial+\Delta$，
则 Jacobi 恒等式给出
\begin{align*}
 [[D,P_\mu],\mathcal O]
 &=[D,\partial_\mu\mathcal O]-[P_\mu,L_D\mathcal O]\\
 &=\partial_\mu(L_D\mathcal O)-L_D(\partial_\mu\mathcal O)
 =\partial_\mu\mathcal O.
\end{align*}
这里第一项通过对 $[D,\mathcal O(x)]$ 求导计算；第二项中坐标系数保持不变。
这确定了 $[D,P_\mu]=P_\mu$ 的号。
同样计算其余交换子，补全为
\begin{align}\label{eq:ghs-conformal-remaining-commutators}
 [M_{\mu\nu},K_\rho]&=\delta_{\nu\rho}K_\mu-\delta_{\mu\rho}K_\nu,\\
 [M_{\mu\nu},M_{\rho\sigma}]&=\delta_{\nu\rho}M_{\mu\sigma}
 -\delta_{\mu\rho}M_{\nu\sigma}
 -\delta_{\nu\sigma}M_{\mu\rho}
 +\delta_{\mu\sigma}M_{\nu\rho},\notag\\
 [D,M_{\mu\nu}]&=0,\qquad [P_\mu,P_\nu]=[K_\mu,K_\nu]=0.\notag
\end{align}
所有这些生成元都是无量纲的表示算符。
若恢复半径 $R$ 的物理圆柱，径向哈密顿量 与 $D$ 的关系是 $H=\hbar cD/R$，
而这里的 $P_\mu$ 不能直接等同于前章具有动量量纲的 $\hat p_\mu$。

\section{初级态、后裔与 Gram 内积}

\begin{definition}[初级态与后裔表示]
设 $|\mathcal O,a\rangle$ 在旋转群下属于某个不可约表示，其中 $a$ 为旋转指标。
初级态满足 $K_\mu|\mathcal O,a\rangle=0$、$D|\mathcal O,a\rangle=\Delta|\mathcal O,a\rangle$。
第 $n$ 级后裔空间由 $P_{\mu_1}\cdots P_{\mu_n}|\mathcal O,a\rangle$ 的线性组合生成。
局域场实现中写 $|\mathcal O\rangle=\mathcal O(0)|0\rangle$，于是
\begin{equation}\label{eq:ghs-primary-descendants}
 D|\mathcal O\rangle=\Delta|\mathcal O\rangle,
 \qquad K_\mu|\mathcal O\rangle=0,
 \qquad D P_{\mu_1}\cdots P_{\mu_n}|\mathcal O\rangle
 =(\Delta+n)P_{\mu_1}\cdots P_{\mu_n}|\mathcal O\rangle.
\end{equation}
最后一个等式逐次使用 $[D,P_\mu]=P_\mu$，说明同一个局域初级算符必然带着整座后裔能级塔。在圆柱上，这些态的能量相差整数倍 $\hbar c/R$。一般相互作用理论仍有此整数能隙结构，但初级维数 $\Delta$ 不必等于自由粒子能量之和。

$D$ 谱有下界，因而初级态是径向能量的最低级态。
把 $K_\mu$ 纳入湮灭初级态的正根子代数时，代数术语把这一构造称为最高权表示；
这不意味着初级态的 $D$ 本征值大于后裔。
\end{definition}

\begin{definition}[径向伴随结构]
在由初级态和后裔组成的稠密代数子空间上，选取对第二变量线性的半正定厄密内积，要求
\begin{equation}\label{eq:ghs-radial-generator-adjoint}
 D^\dagger=D,\qquad P_\mu^\dagger=K_\mu,
 \qquad M_{\mu\nu}^\dagger=-M_{\mu\nu}.
\end{equation}

令 $\mathcal N$ 是全部零范数方向，取商空间后再完备化，才得到正定 Hilbert 空间。
这些伴随关系是这里研究的表示结构的一部分，不是对坐标微分表达式逐项取共轭的结果。
在具有反射正性的共形场论中，球外区域反演到球内把径向时间反向，
并将平移变成特殊共形变换，因而实现此伴随结构；该物理构造需额外的态制备与正性假设。
\end{definition}

\begin{definition}[后裔 Gram 矩阵]
在第 $n$ 级选择线性生成组 $|n,A\rangle$，定义
\[
 G^{(n)}_{AB}=\langle n,A|n,B\rangle.
\]
任意系数向量 $c$ 的范数是 $c^\dagger G^{(n)}c$，所以幺正性要求所有级的 Gram 矩阵半正定。
线性相关生成元和真正的零范数后裔都要取商；在取商之前不能直接反演 Gram 矩阵。
利用 $P^\dagger=K$，每一矩阵元都可通过交换 $K$ 与 $P$、最后使用 $K|\mathcal O\rangle=0$ 来计算。
\end{definition}

\begin{example}[归一标量的一级 Gram 矩阵]
给定归一化标量初级态，用 $K_\mu=P_\mu^\dagger$ 计算一级后裔的 Gram 矩阵，并由正性判断维数的限制。

\textbf{解。} 令 $\langle\mathcal O|\mathcal O\rangle=1$ 且 $M_{\mu\nu}|\mathcal O\rangle=0$，则
\begin{align}\label{eq:ghs-scalar-level-one-gram}
 \langle P_\mu\mathcal O|P_\nu\mathcal O\rangle
 &=\langle\mathcal O|K_\mu P_\nu|\mathcal O\rangle\\
 &=\langle\mathcal O|[K_\mu,P_\nu]|\mathcal O\rangle
 =2\Delta\delta_{\mu\nu}.\notag
\end{align}
它的非负性至少要求 $\Delta\ge0$。然而对于非平凡标量，这还不是最强的界；必须再考察二级迹后裔。
\end{example}

\begin{figure}[htbp]
\centering
% Local TikZ geometry; styles inherited exclusively from the shared preamble.
\begin{tikzfig}[scale=1]
\node[draw,circle,minimum size=6mm] (p) at (0,0) {$\mathcal O$};
\node[draw,circle,minimum size=6mm] (a) at (2.55,0) {$P_\mu\mathcal O$};
\node[draw,circle,minimum size=6mm] (b) at (5.3,0) {$P_{(\mu}P_{\nu)}\mathcal O$};
\draw[lecturearrow] (p) to[bend left=21] node[above] {$P_\mu$} (a);
\draw[lecturearrow] (a) to[bend left=21] node[above] {$P_\nu$} (b);
\draw[lecturearrow] (a) to[bend left=21] node[below] {$K_\mu$} (p);
\draw[lecturearrow] (b) to[bend left=21] node[below] {$K_\nu$} (a);
\node at (0,-1.06) {Level 0};
\node at (2.55,-1.06) {Level 1};
\node at (5.3,-1.06) {Level 2};
\node at (0,.98) {$\Delta$};\node at (2.55,.98) {$\Delta+1$};\node at (5.3,.98) {$\Delta+2$};
\end{tikzfig}

\caption{共形初级态和前两级后裔的标度维数。$P_\mu$ 提升级数，$K_\mu$ 降低级数；图仅展示部分张量分量，同一级的真实 Gram 空间一般具有多个独立方向。}\label{fig:ghs-descendant-weight-levels}
\end{figure}

\cref{fig:ghs-descendant-weight-levels} 中二级后裔有迹与对称无迹分量，因而幺正性界必须通过 Gram 矩阵的本征值而不只是一级范数来判定。

\section{幺正性界与表示缩短}

\begin{theorem}[标量的必要幺正性界]
设 $d>2$，标量初级态非零且 $\Delta>0$。若其后裔表示具有上述半正定内积，则
\begin{equation}\label{eq:ghs-scalar-unitarity-bound}
 \Delta\ge\frac{d-2}{2}.
\end{equation}

若 $\Delta=(d-2)/2$，二级迹后裔 $P^2|\mathcal O\rangle$ 是零范数态，必须在正定表示中模掉。
\end{theorem}
\begin{proof}
对一级后裔再作用平移，会得到 $P_\mu P_\nu|\mathcal O\rangle$。将迹与对称无迹部分分开，迹态为 $P^2|\mathcal O\rangle=\sum_\mu P_\mu P_\mu|\mathcal O\rangle$。它的范数可以完全用代数计算。


 由 $[K_\mu,P_\nu]=2\delta_{\mu\nu}D-2M_{\mu\nu}$，对标量初级态有
 \begin{align*}
 K_\mu P^2|\mathcal O\rangle
 &=[K_\mu,P_\nu]P_\nu|\mathcal O\rangle
   +P_\nu[K_\mu,P_\nu]|\mathcal O\rangle\\
 &=2D P_\mu|\mathcal O\rangle
   -2M_{\mu\nu}P_\nu|\mathcal O\rangle
   +2P_\mu D|\mathcal O\rangle.
 \end{align*}
 在此约定下 $[M_{\mu\nu},P_\nu]=(d-1)P_\mu$，所以 $M_{\mu\nu}P_\nu|\mathcal O\rangle=(d-1)P_\mu|\mathcal O\rangle$。又因为 $DP_\mu|\mathcal O\rangle=(\Delta+1)P_\mu|\mathcal O\rangle$，故
 \begin{align*}
 K_\mu P^2|\mathcal O\rangle
 &=\bigl[2(\Delta+1)-2(d-1)+2\Delta\bigr]P_\mu|\mathcal O\rangle\\
 &=4\left(\Delta-\frac{d-2}{2}\right)P_\mu|\mathcal O\rangle.
 \end{align*}
 再左乘 $\langle\mathcal O|K_\mu$ 并对 $\mu$ 求和，利用\cref{eq:ghs-scalar-level-one-gram} 得
 \begin{equation*}
 \|P^2|\mathcal O\rangle\|^2
 =\sum_\mu\langle\mathcal O|K_\mu K_\mu P^2|\mathcal O\rangle
 =8d\Delta\left(\Delta-\frac{d-2}{2}\right).
 \end{equation*}

因为 $d>2$ 且 $\Delta>0$，最后一式非负恰要求\eqnrefs{eq:ghs-scalar-unitarity-bound}。
在饱和点，该迹态范数为零；$\Delta=0$ 的标量不在本命题的假设内。
\end{proof}

正定商空间中的零向量对所有态都正交，这由半正定内积的 Cauchy--Schwarz 不等式推出。
因此在局域算符实现中，迹后裔脱耦写成 $\partial^2\mathcal O=0$。
它给出饱和标量的自由波动方程，也解释了为何运动方程能消掉看似独立的导数算符。
这一推断使用局域算符与态的对应；抽象群表示本身只陈述 $P^2|\mathcal O\rangle=0$。

\begin{theorem}[对称无迹自旋的必要幺正性界]
设 $J_{\mu_1\ldots\mu_S}$ 为非零秩 $S\geq1$ 的对称无迹初级态，初级旋转多重态内积正定。
若其全部后裔的内积半正定，则
\begin{equation}\label{eq:ghs-spin-unitarity-bound-extended}
 \Delta\ge S+d-2,\qquad S\ge1.
\end{equation}

饱和时，散度后裔 $P^\nu|J_{\nu\mu_2\ldots\mu_S}\rangle$ 在正定商空间中为零。
\end{theorem}
\begin{proof}
 把对称无迹初级态记为 $|J_{\mu_1\cdots\mu_S}\rangle$。其标度维数是 $\Delta$，$K_\alpha|J\rangle=0$。旋转生成元在每个指标上分别作用：
 \begin{align*}
 M_{\alpha\nu}|J_{\rho_1\cdots\rho_S}\rangle
 =\sum_{j=1}^S\bigl(
 \delta_{\nu\rho_j}|J_{\rho_1\cdots\alpha\cdots\rho_S}\rangle
 -\delta_{\alpha\rho_j}|J_{\rho_1\cdots\nu\cdots\rho_S}\rangle\bigr).
 \end{align*}
 定义散度后裔 $|L_{\mu_2\cdots\mu_S}\rangle
 =\sum_\nu P_\nu|J_{\nu\mu_2\cdots\mu_S}\rangle$。因为 $K_\alpha$ 湮灭初级态，
 \begin{align*}
 K_\alpha|L_{\mu_2\cdots\mu_S}\rangle
 &=\sum_\nu[K_\alpha,P_\nu]
 |J_{\nu\mu_2\cdots\mu_S}\rangle\\
 &=2\Delta|J_{\alpha\mu_2\cdots\mu_S}\rangle
 -2\sum_\nu M_{\alpha\nu}|J_{\nu\mu_2\cdots\mu_S}\rangle.
 \end{align*}
 现在按旋转作用的指标逐项求和。第一个指标贡献
 \begin{equation*}
 \sum_\nu\bigl(\delta_{\nu\nu}J_{\alpha\mu_2\cdots}
 -\delta_{\alpha\nu}J_{\nu\mu_2\cdots}\bigr)
 =(d-1)J_{\alpha\mu_2\cdots}.
 \end{equation*}
 其余每个指标 $\mu_j$ 的第一项，借对称性各贡献一个 $J_{\alpha\mu_2\cdots}$；第二项含 $\sum_\nu J_{\nu\nu\cdots}=0$，因无迹性而消失。因此
 \begin{equation*}
 \sum_\nu M_{\alpha\nu}|J_{\nu\mu_2\cdots\mu_S}\rangle
 =(d+S-2)|J_{\alpha\mu_2\cdots\mu_S}\rangle.
 \end{equation*}
 合并前式后得到完整的散度后裔关系
 \begin{equation*}
 K_\alpha|L_{\mu_2\cdots\mu_S}\rangle
 =2(\Delta-d-S+2)|J_{\alpha\mu_2\cdots\mu_S}\rangle.
 \end{equation*}
 最后用 $P_\nu^\dagger=K_\nu$ 计算所有张量分量范数之和：
 \begin{align*}
 \sum_{\mu_2,\ldots,\mu_S}\|L_{\mu_2\cdots\mu_S}\|^2
 &=\sum_{\nu,\mu_2,\ldots,\mu_S}
 \langle J_{\nu\mu_2\cdots\mu_S}|K_\nu
 |L_{\mu_2\cdots\mu_S}\rangle\\
 &=2(\Delta-d-S+2)
 \sum_{\nu,\mu_2,\ldots,\mu_S}
 \|J_{\nu\mu_2\cdots\mu_S}\|^2.
 \end{align*}
 右侧最后一个求和严格为正，除非初级态为零。因此 $\Delta\ge d+S-2$；等号成立时，散度后裔整体为零范数。这个证明对任意整数 $S\ge1$ 都适用。
\end{proof}

在局域共形场论中，饱和关系给出 $\partial^{\mu_1}J_{\mu_1\ldots\mu_S}=0$。
证明只用了共形代数、旋转张量作用和正性，没有使用 Wick 配对或自由场方程。
因此即使无法把初级算符写成基本场的显式乘积，界与缩短关系仍然成立。
这些界是表示存在的必要条件；一个满足界的数值并不单独构造出完整的局域场论。

\begin{example}[矢量散度通道的直接核对]
对归一化矢量初级态，直接计算散度后裔的范数，求该通道给出的幺正性条件。

\textbf{解。} 令 $|J_\mu\rangle$ 是归一化矢量初级态，$\langle J_\mu|J_\nu\rangle=\delta_{\mu\nu}$，并记 $|\partial\cdot J\rangle=\sum_\nu P_\nu|J_\nu\rangle$。矢量表示中的旋转作用为
 \begin{equation*}
 M_{\mu\nu}|J_\rho\rangle
 =\delta_{\nu\rho}|J_\mu\rangle-\delta_{\mu\rho}|J_\nu\rangle.
 \end{equation*}
 将它代入交换子，逐项求和：
 \begin{align*}
 \|\partial\cdot J\|^2
 &=\sum_{\mu,\nu}\langle J_\mu|K_\mu P_\nu|J_\nu\rangle\\
 &=2\sum_{\mu,\nu}\left[\delta_{\mu\nu}\Delta
 \langle J_\mu|J_\nu\rangle
 -\langle J_\mu|M_{\mu\nu}|J_\nu\rangle\right]\\
 &=2d\Delta-2d(d-1)=2d(\Delta-d+1).
 \end{align*}
 从而 $\Delta\ge d-1$，恰好是 \cref{eq:ghs-spin-unitarity-bound-extended} 在 $S=1$ 的情形。
\end{example}

这一推导还表明：所谓“守恒”并不是任意把散度设成零，而是表示论中某个后裔与物理 Hilbert 空间脱耦。自由标量的双线性流满足 $\Delta_S=S+d-2$，因而恰好处在这一边界；一般相互作用使 $\Delta_S$ 增大时，原来被模掉的后裔重新取得正范数。


\begin{proposition}[零范数子空间与短表示]
若所有后裔级的 $G^{(n)}$ 均半正定，则全部零范数态在 $P,K,D,M$ 作用下构成不变子空间。
把这一子空间模掉得到短表示。对于饱和的标量或对称无迹张量，
运动方程或守恒方程消去的并非任意挑选的分量，而是上面已经算出的 Gram 零方向及其后裔。
\end{proposition}
\begin{proof}
若 $v$ 范数为零，半正定 Cauchy--Schwarz 不等式给出 $\langle w,v\rangle=0$ 对所有代数态 $w$ 成立。
对任意生成元 $T$，其伴随仍是同一生成元集合中的算符，故
\[
 \langle w,Tv\rangle=\langle T^\dagger w,v\rangle=0.
\]
取 $w=Tv$ 得到 $Tv$ 也为零范数。所有生成元都保持零空间，商表示因而良定义。
\end{proof}

若对称无迹初级态的维数从 $S+d-2$ 增为 $S+d-2+\gamma_S$，
散度范数变为 $2\gamma_S$ 乘以原初级多重态的范数和。
因此后裔重新取得非零范数，对应短表示向长表示的变化；要由此计算 $\gamma_S$，
还须指定相互作用、算符混合及散度算符的实际二点归一化。
群论至此给出计算所需的代数约束，场论再决定这些动力学数据。
